\begin{proof}[Beweis des Lemmas der Existenz einer tubularen Umgebung\label{jordan:proof:tubulare_umgebung}]
    Sei $\gamma$ eine reell-analytische Kurve in der Ebene $\mathbb{R}^2$ und $\varphi$ die durch $\varphi(t,s) = \gamma(t) + s \cdot n(t)$ definierte Abbildung. 
    Die Jacobi-Determinante von $\varphi$ in $(t,0)$ ist $||\gamma'(t)|| \neq 0$. 
    Nach dem Satz über die Umkehrfunktion ist $\varphi$ lokal ein Diffeomorphismus um jeden Punkt $(t,0)$. 
    Globale Injektivität auf einem Streifen $S\textsuperscript{1} \times (-\epsilon, \epsilon)$ folgt aus der Kompaktheit und der Einfachheit von $\gamma$: 
    Wäre $\varphi(t_1,s_1) = \varphi(t_2,s_2)$ für $(t_1,s_1) \neq (t_2,s_2)$, so würde dies bedeuten, dass $\gamma(t_1) + s_1 n(t_1) = \gamma(t_2) + s_2 n(t_2)$, was im Widerspruch zur Injektivität von $\gamma$ steht.
\end{proof}